\chapter{Limits: Intuition and the Epsilon-Delta Definition}

In the previous chapter you met the \emph{idea} of a limit. You watched values
of a function pile up near a single number $L$ as the input crept toward some
point $a$, and you learned to write
\[
  \lim_{x \to a} f(x) = L
\]
to record that fact. You also saw that the early calculus -- Newton's
``ultimate ratios,'' Leibniz's infinitesimals -- ran on exactly this kind of
intuition, and that for a century and a half nobody could say with full
precision what the words actually \emph{meant}. ``Approaches.'' ``Tends to.''
``Becomes infinitely close.'' These are pictures, not definitions. A picture
cannot be argued with; a definition can.

This chapter makes the informal concept of a limit rigorous. We take the intuitive
phrase ``$f(x)$ gets close to $L$ as $x$ gets close to $a$'' and formulate it
as a precise mathematical statement: the \emph{epsilon-delta definition}. Then we use
that statement to prove results directly: verifying specific limits, and establishing
that the familiar limit laws (such as the limit of a sum being the sum of the limits)
follow deductively rather than standing as unproven assumptions.

Written in symbols, the epsilon-delta definition can look dense because it stacks ``for every'' ($\forall$) and ``there exists'' ($\exists$) in a single sentence. Underneath the symbols, though, it is just the game of \emph{challenge and response} we previewed at the end of Chapter~5: a skeptic names an output tolerance, and we find an input window narrow enough to meet it. Let us play that game in plain English first, and let the symbols grow out of it.

\chapterline

\section{Why intuition alone is not enough}
\label{ch06:why-not-enough}

You might reasonably ask: if the intuitive picture works, if everyone
\emph{agrees} that $\lim_{x\to 2}(3x+1) = 7$, why bother making it precise? The
answer is that intuition, left to itself, contradicts itself. It gives
different answers to different people, and sometimes the wrong answer to
everyone.

\begin{intuitionbox}
Consider the sloppy slogan ``the limit is the value the function gets close to.''
Now look at
\[
  f(x) = \frac{x^2 - 1}{x - 1}.
\]
At $x = 1$ this is $\tfrac{0}{0}$, undefined. But for every $x \ne 1$ we can
cancel: $f(x) = x + 1$. So as $x$ heads toward $1$, $f(x)$ heads toward $2$. The
function never \emph{reaches} $2$ at the relevant input -- it is not even
defined at $x=1$ -- yet $2$ is unmistakably the limit. ``Gets close to'' has to
mean something that does not depend on the value at $a$ at all. The intuitive
phrase does not say that out loud, so different readers fill the gap
differently. Figure~\ref{ch06:fig-hole} shows the situation: a removable
discontinuity -- a single missing point or hole punched out of an otherwise
unbroken line.
\end{intuitionbox}

\begin{figure}[htbp]
\centering
\begin{tikzpicture}
\begin{axis}[
  width=8.4cm, height=5.6cm,
  title style={font=\small},
  title={Removable discontinuity: approaching the hole},
  axis lines=middle, xlabel={$x$}, ylabel={$y$},
  xtick={1}, xticklabels={$1$}, ytick={2}, yticklabels={$2$},
  xmin=-0.6, xmax=2.6, ymin=-0.4, ymax=3.6, clip=false,
]
\addplot[curveblue, very thick, domain=-0.6:2.6, samples=2] {x+1};
\node[curveblue, anchor=south east] at (axis cs:2.1,3.18) {$y=x+1$};
\draw[gray, dashed] (axis cs:1,-0.4) -- (axis cs:1,2);
\draw[gray, dashed] (axis cs:-0.6,2) -- (axis cs:1,2);
\fill[white] (axis cs:1,2) circle (2.6pt);
\draw[curveblue, very thick] (axis cs:1,2) circle (2.6pt);
\end{axis}
\end{tikzpicture}
\caption{The graph of $f(x)=\frac{x^2-1}{x-1}$. For every $x\ne1$ it agrees
exactly with the line $y=x+1$; only at $x=1$ is it undefined, leaving a
removable discontinuity (an isolated hole). The limit asks where the graph is
\emph{heading} as $x\to1$, which is plainly $2$ -- the very spot the hole sits
over. The missing point is exactly what the limit is built to ignore.}
\label{ch06:fig-hole}
\end{figure}

Worse, the eighteenth-century notion of an ``infinitely small'' quantity $\dd x$
-- a number bigger than zero but smaller than every ordinary positive number --
is, taken literally, self-contradictory in the ordinary real numbers. If
$\dd x > 0$ and $\dd x$ is smaller than every positive number, then $\dd x$ is
smaller than $\dd x / 2$, which is absurd. The early users got correct answers
\emph{despite} this, by a kind of trained instinct for when the infinitesimals
were safe to drop. What they lacked was a definition that a skeptic could check
mechanically, without instinct.

\begin{historybox}
The skeptic arrived. In 1734 Bishop George Berkeley published \emph{The Analyst},
a sharp critique of the foundations of calculus, asking how a quantity $\dd x$
could be treated as nonzero (so one can divide by it) and then as zero (so one
can discard the remainder) within the same derivation. Berkeley was not being
obtuse; he was right that no one could yet answer him. It took until the 1820s
for Augustin-Louis Cauchy to base calculus on limits rather than infinitesimals,
and until Karl Weierstrass's lectures in the 1850s--1870s for the limit itself
to be pinned down in the static, quantifier-based form -- epsilon and delta --
that we use today. The definition in this chapter is, almost word for word,
Weierstrass's answer to Berkeley's critique.
\end{historybox}

So the goal is a definition with three properties: it must (1) not refer to the
value $f(a)$ at all, (2) capture ``arbitrarily close'' without infinitesimal
numbers, using only ordinary positive real numbers, and (3) be
checkable -- a finite, rigorous criterion that settles whether a proposed $L$ is
the limit. A single definition accomplishes all three.

\section{The challenge--response game}
\label{ch06:game}

Let me describe the game first, with no symbols, and then we will name the
pieces.

Two players sit across a table. Player One is the \emph{Skeptic}; Player Two is
the \emph{Prover}. On the table is a claim: ``As $x \to a$, $f(x) \to L$.'' The
Skeptic does not believe it and wants to test it. Here is the single move each
player gets, repeated as often as the Skeptic likes.

\begin{focusbox}
\textbf{The move.} The Skeptic names a \emph{tolerance}: a positive number, a
margin of error around $L$ that they are willing to accept. Call it $\eps$
(epsilon, for ``error''). The Skeptic says: ``Fine. I want $f(x)$ to land
within $\eps$ of $L$. Can you guarantee that?''

The Prover must respond with a \emph{closeness}: a positive number $\delta$
(delta, for ``distance''), a radius around the input $a$, and a promise: ``Yes.
As long as you keep $x$ within $\delta$ of $a$ (but not exactly at $a$), I
guarantee $f(x)$ stays within your $\eps$ of $L$.''
\end{focusbox}

The Prover \emph{wins a round} if the promise holds: every $x$ in the punctured
interval $(a-\delta,\,a+\delta)\setminus\{a\}$ -- an open interval around $a$
with the center point $a$ itself removed -- really does produce an $f(x)$
inside $(L-\eps,\,L+\eps)$. The Prover \emph{wins the game} -- and the limit
really is $L$ -- if the Prover has a winning response \emph{no matter what
tolerance the Skeptic names}, including absurdly tiny ones like $\eps =
10^{-100}$.

Pause on the asymmetry, because it is the whole point.

\begin{intuitionbox}
The Skeptic moves \emph{first}. The Skeptic gets to pick $\eps$ as small as they
please, and only \emph{then} does the Prover get to choose $\delta$. This is
exactly right. ``$f(x)$ gets arbitrarily close to $L$'' should mean: name any
closeness you want for the \emph{output}, and I can deliver it by squeezing the
\emph{input} hard enough. The Prover's $\delta$ is allowed to depend on the
Skeptic's $\eps$. Usually it must: the smaller the output tolerance, the tighter
the Prover has to clamp the input. A winning strategy is a \emph{rule} that
turns any $\eps$ into a $\delta$ that works.
\end{intuitionbox}

If even one stubborn tolerance defeats the Prover -- one $\eps$ for which no
$\delta$ keeps the function inside the band -- then the limit is not $L$. That is
how the game lets us \emph{disprove} a wrong claim too, which we will use later.

\begin{walkbox}
Before reading on, play a round in your head with $f(x) = 3x + 1$ near $a = 2$,
claiming $L = 7$. The Skeptic says $\eps = 0.6$: ``land within $0.6$ of $7$, so
between $6.4$ and $7.6$.'' What inputs $x$ near $2$ give outputs in that band?
Since $3x+1$ is between $6.4$ and $7.6$ exactly when $x$ is between $1.8$ and
$2.2$, the Prover can answer $\delta = 0.2$. Now ask yourself the move that
matters: if the Skeptic had said $\eps = 0.06$ instead, what $\delta$ would the
Prover answer? Notice the pattern in how $\delta$ shrinks with $\eps$. That
pattern \emph{is} the proof.
\end{walkbox}

\subsection{Two pictures: a band and an interval}

The game lives in a picture, and you should carry the picture with you forever.
The Skeptic's $\eps$ is a horizontal \emph{band} of height $2\eps$ centered on
the line $y = L$. The Prover's $\delta$ is a vertical \emph{strip} of width
$2\delta$ centered on $x = a$. The Prover wins the round when the graph of $f$,
restricted to that vertical strip, stays trapped inside the horizontal band --
with the single point above $x=a$ exempted, since we never ask about $x=a$
itself. See Figure~\ref{ch06:fig-band}.

\begin{figure}[htbp]
\centering
\begin{tikzpicture}
\begin{axis}[
  width=12cm, height=7.5cm,
  title style={font=\small},
  title={The challenge--response game: horizontal error band vs. vertical input strip},
  axis lines=middle,
  xlabel={$x$}, ylabel={$y$},
  xtick={1.8, 2, 2.2}, xticklabels={$a-\delta$, $a$, $a+\delta$},
  ytick={5.8, 7, 8.2}, yticklabels={$L-\eps$, $L$, $L+\eps$},
  xmin=0.6, xmax=3.4, ymin=2.5, ymax=11,
  clip=false,
]
\addplot[draw=none, fill=green!12] coordinates {(0.6,8.2)(3.4,8.2)(3.4,5.8)(0.6,5.8)};
\addplot[draw=none, fill=blue!10] coordinates {(1.8,2.5)(2.2,2.5)(2.2,11)(1.8,11)};
\draw[green!50!black, dashed] (axis cs:0.6,8.2) -- (axis cs:3.4,8.2);
\draw[green!50!black, dashed] (axis cs:0.6,5.8) -- (axis cs:3.4,5.8);
\draw[blue!50!black, dashed] (axis cs:1.8,2.5) -- (axis cs:1.8,11);
\draw[blue!50!black, dashed] (axis cs:2.2,2.5) -- (axis cs:2.2,11);
\addplot[curveblue, very thick, domain=0.8:3.2, samples=2] {3*x+1};
\node[curveblue, anchor=south east] at (axis cs:2.9,9.8) {$y=f(x)$};
\draw[gray, dashed] (axis cs:2,2.5) -- (axis cs:2,7);
\draw[gray, dashed] (axis cs:0.6,7) -- (axis cs:2,7);
\fill[white] (axis cs:2,7) circle (2.4pt);
\draw[curveblue] (axis cs:2,7) circle (2.4pt);
\end{axis}
\end{tikzpicture}
\caption{The challenge--response game as a picture. The Skeptic paints the
horizontal band of half-height $\eps$ about $y=L$. The Prover must paint a
vertical strip of half-width $\delta$ about $x=a$ \emph{narrow enough} that the
graph, inside the strip, never escapes the band. The open circle marks that
$x=a$ is excluded: we care about the approach, not the value there.}
\label{ch06:fig-band}
\end{figure}

\begin{happeningbox}
What is the limit definition really asserting? Not that ``$f(x)$ moves toward
$L$'' -- the notion of ``moving'' suggests time and motion, which the static real numbers do
not contain. Instead, the definition asserts a condition on \emph{containment of sets}.
For every horizontal band around $L$, however thin, there exists a vertical strip
around $a$, possibly very thin, such that the part of the graph over the
punctured strip lies entirely inside the band. No motion, no infinitely small
quantities, and no reference to $f(a)$ are needed. Just: ``for every band, there is a strip that
fits inside it.'' Weierstrass's insight was that this containment condition
captures the intended meaning of ``approaches'' in terms that can be verified precisely.
\end{happeningbox}

\section{From the game to the definition}
\label{ch06:translate}

Now we translate, phrase by phrase. This is the moment to go slowly; if you
internalize the translation, the definition will never again look foreign.

\begin{center}
\renewcommand{\arraystretch}{1.5}
\begin{tabular}{@{}p{6.3cm}@{\quad$\Longleftrightarrow$\quad}p{6.1cm}@{}}
\toprule
\textbf{Game language} & \textbf{Symbol language} \\
\midrule
``$f(x)$ lands within $\eps$ of $L$'' & $\abs{f(x) - L} < \eps$ \\
``$x$ is within $\delta$ of $a$, but not at $a$'' & $0 < \abs{x - a} < \delta$ \\
``the Skeptic may name any tolerance'' & for every $\eps > 0$ \\
``the Prover can find a closeness'' & there exists $\delta > 0$ \\
``the promise holds for those $x$'' & ($0<\abs{x-a}<\delta$) implies ($\abs{f(x)-L}<\eps$) \\
\bottomrule
\end{tabular}
\end{center}

Two of these rows deserve a second look.

First, why $\abs{f(x)-L} < \eps$ and not $\le \eps$? Because ``within
$\eps$'' is a strict band, and -- more importantly -- it does not matter. Any
tolerance the Skeptic can hit with $\le$ they can also hit with $<$ by switching
to a slightly larger $\eps$, so the two versions of the definition are
equivalent. We use strict inequalities throughout for clean algebra.

Second, and this is the subtle one: why the \emph{left} inequality
$0 < \abs{x-a}$? The expression $\abs{x-a}$ is zero precisely when $x = a$.
Writing $0 < \abs{x-a}$ forbids that single case. This is the formal incarnation
of ``we never ask about $x = a$ itself.'' It is what makes the limit blind to
$f(a)$ -- exactly the property we demanded in
Section~\ref{ch06:why-not-enough}, the one that let
$\frac{x^2-1}{x-1}$ have a limit of $2$ at a point where it is undefined.

Stack the rows and you have the definition. Read it now not as a wall of
symbols but as the game you just played.

\begin{definition}[The limit, Weierstrass]
\label{ch06:def-limit}
Let $f$ be defined on some open interval containing $a$, except possibly at $a$
itself. We say
\[
  \lim_{x \to a} f(x) = L
\]
if for every $\eps > 0$ there exists a $\delta > 0$ such that, for all $x$,
\[
  0 < \abs{x - a} < \delta
  \quad\Longrightarrow\quad
  \abs{f(x) - L} < \eps.
\]
\end{definition}

\begin{formalbox}
Read the quantifiers in order, left to right, as alternating players:
\[
  \underbrace{\forall\, \eps > 0}_{\text{Skeptic}}\;\;
  \underbrace{\exists\, \delta > 0}_{\text{Prover}}\;\;
  \underbrace{\forall x\,\big(0<\abs{x-a}<\delta \Rightarrow \abs{f(x)-L}<\eps\big)}_{\text{the promise is kept}}.
\]
The \emph{order} of quantifiers is essential. Because $\delta$ comes after $\eps$,
$\delta$ may be chosen in terms of $\eps$---it typically depends directly on $\eps$.
If the order were reversed to ``$\exists\delta\,\forall\eps$,''
one would be demanding a single universal $\delta$ that works for all tolerances simultaneously,
which is impossible for non-constant functions. The quantifier ordering encodes
the dependency.
\end{formalbox}

\begin{trapbox}
\textbf{The value at $a$ is irrelevant -- truly irrelevant.} Three different
functions can share the same limit at $a$: one undefined there, one defined with
the ``right'' value $L$, and one defined with a deliberately wrong value. Because
Definition~\ref{ch06:def-limit} only ever quantifies over $x$ with
$\abs{x-a}>0$, none of these differences at the single point $a$ can change the
limit. Beginners routinely ``simplify'' by plugging in $x=a$; that is a
\emph{shortcut} that happens to work for nice (continuous) functions, not the
definition. We will see in the next chapter exactly when the shortcut is legal.
\end{trapbox}

\subsection{What it would take to be \emph{false}}

A definition you cannot negate is a definition you do not understand. The limit
fails to be $L$ when the Skeptic has a winning tolerance. Negating the
quantifiers (each $\forall$ becomes $\exists$ and vice versa) gives:

\begin{focusbox}
$\displaystyle \lim_{x\to a} f(x) \ne L$ means: \emph{there exists} some
$\eps > 0$ such that \emph{for every} $\delta > 0$ there is at least one $x$ with
$0 < \abs{x-a} < \delta$ but $\abs{f(x)-L} \ge \eps$.

In words: there is one bad band the Prover can never fit into, no matter how
thin a strip they try -- every strip leaks at least one point of the graph out
of that band.
\end{focusbox}

\begin{recoverybox}
\textbf{Pause and think.} Suppose someone claims that $\lim_{x\to 0} f(x) = 5$.
To disprove their claim using the negated definition, what exactly do you need
to produce? You do not need to check every $\eps$. You only need to name
\emph{one} stubborn error tolerance $\eps_0 > 0$ such that every interval around
$0$, no matter how tiny, contains an input $x$ whose output $f(x)$ stays at
least $\eps_0$ away from $5$. A single leaky band defeats the Prover forever.
\end{recoverybox}

We will put this to work in Section~\ref{ch06:no-limit} to show that
$\lim_{x\to 0}\sin(1/x)$ does not exist: there is a fixed band that every strip,
however thin, fails to contain because the graph oscillates forever.

\section{Worked epsilon--delta proofs: linear functions}
\label{ch06:linear}

Time to prove something. The structure of every epsilon-delta proof is the same,
and once you see it twice it becomes a template:

\begin{enumerate}
\item \textbf{Scratch work (find $\delta$).} Start from the goal
  $\abs{f(x)-L}<\eps$ and work \emph{backwards}, unraveling it into a condition
  on $\abs{x-a}$. This tells you what $\delta$ to pick. This step is
  exploration; it does not go in the final proof.
\item \textbf{The proof (verify $\delta$ works).} Announce the $\delta$ you
  found. Then, starting from $0<\abs{x-a}<\delta$, run the algebra
  \emph{forward} to reach $\abs{f(x)-L}<\eps$.
\end{enumerate}

The scratch work is where the thinking happens; the proof is where you check it.
Let me show both halves out loud the first time, then tighten up.

\begin{example}[A linear limit]
\label{ch06:ex-linear}
Prove that $\displaystyle \lim_{x\to 2}(3x+1) = 7$.
\end{example}

\textbf{Scratch work.} The output gap is
\[
  \abs{f(x) - L} = \abs{(3x+1) - 7} = \abs{3x - 6} = 3\abs{x-2}.
\]
There it is: the output distance is exactly $3$ times the input distance. We want
$3\abs{x-2} < \eps$, i.e. $\abs{x-2} < \eps/3$. So if the Skeptic hands us
$\eps$, we should respond $\delta = \eps/3$. (Now we know the answer; the proof
below pretends we are seeing it for the first time and just verifies it.)

\textbf{Proof.} Let $\eps > 0$ be given. Choose $\delta = \eps/3$. Suppose
$0 < \abs{x-2} < \delta$. Then
\[
  \abs{(3x+1) - 7} = 3\abs{x-2} < 3\delta = 3\cdot\frac{\eps}{3} = \eps.
\]
Thus $0<\abs{x-2}<\delta$ implies $\abs{(3x+1)-7}<\eps$. Since $\eps>0$ was
arbitrary, $\lim_{x\to 2}(3x+1) = 7$. $\qquad\blacksquare$

Notice how short the actual proof is, and how completely the scratch work did the
work. The Skeptic-to-Prover rule here is the function $\delta = \eps/3$: hand it
any tolerance and it returns a closeness that provably works.

\begin{whybox}
Why does $\delta = \eps/3$ work, and why \emph{that} number? Because a line of
slope $3$ magnifies horizontal distances into vertical ones by exactly the factor
$3$. If you want the vertical wandering held under $\eps$, you must hold the
horizontal wandering under $\eps/3$ -- you ``pre-divide'' by the magnification.
For a general line $f(x)=mx+b$ with $m\ne 0$, the same reasoning gives
$\delta = \eps/\abs{m}$: the steeper the line, the tighter you must clamp the
input. This is your first glimpse of a theme that runs through all of calculus --
\emph{the local rate of change controls how output error responds to input
error}. The number $\abs{m}$ is, of course, the derivative of the line. We are
already computing with derivatives; we just have not named them yet.
\end{whybox}

\begin{labbox}
Carry out the whole template yourself for $\displaystyle\lim_{x\to -1}(5 - 2x) =
7$. (Scratch: $\abs{(5-2x)-7} = \abs{-2x-2} = 2\abs{x+1}$, so take
$\delta = \eps/2$.) Then do $\lim_{x\to 3}(7) = 7$ for the constant function and
watch something curious happen: the output gap is $\abs{7-7}=0<\eps$
\emph{automatically}, so \emph{any} $\delta$ works -- say $\delta = 1$. The
Prover never even has to look at $\eps$. A constant function is the one case
where the Prover can answer before the Skeptic finishes speaking.
\end{labbox}

\begin{figure}[htbp]
\centering
\begin{tikzpicture}
\begin{axis}[
  width=8.4cm, height=6.4cm,
  title style={font=\small},
  title={Nested tolerances: shrinking $\eps$ forces shrinking $\delta$},
  axis lines=middle, xlabel={$x$}, ylabel={$y$},
  xtick={2}, xticklabels={$2$}, ytick={7}, yticklabels={$7$},
  xmin=1.3, xmax=2.7, ymin=4, ymax=10, clip=false,
]
\addplot[draw=none, fill=green!18] coordinates {(1.3,7.6)(2.7,7.6)(2.7,6.4)(1.3,6.4)};
\addplot[draw=none, fill=green!32] coordinates {(1.3,7.3)(2.7,7.3)(2.7,6.7)(1.3,6.7)};
\addplot[draw=none, fill=blue!12] coordinates {(1.8,4)(2.2,4)(2.2,10)(1.8,10)};
\addplot[draw=none, fill=blue!24] coordinates {(1.9,4)(2.1,4)(2.1,10)(1.9,10)};
\addplot[curveblue, very thick, domain=1.3:2.7, samples=2] {3*x+1};
\node[curveblue, anchor=south east] at (axis cs:2.55,9.0) {$y=3x+1$};
\node[green!45!black, font=\scriptsize, anchor=south west] at (axis cs:1.35,7.62) {$\eps_1=0.6$};
\node[green!45!black, font=\scriptsize, anchor=south west] at (axis cs:1.35,7.32) {$\eps_2=0.3$};
\node[blue!45!black, font=\scriptsize, anchor=south] at (axis cs:2.0,4.2) {$\delta_2=0.1$};
\node[blue!45!black, font=\scriptsize, anchor=south] at (axis cs:1.68,4.6) {$\delta_1=0.2$};
\end{axis}
\end{tikzpicture}
\caption{For the line $3x+1$ at $a=2$: shrink the green band (smaller $\eps$) and
the matching blue strip shrinks in lockstep, by the factor $\delta=\eps/3$. The
nested bands and strips make the rule $\eps\mapsto\eps/3$ visible.}
\label{ch06:fig-linear}
\end{figure}

\section{Worked epsilon--delta proofs: a quadratic}
\label{ch06:quadratic}

Lines are easy because the output gap factors cleanly as a constant times the
input gap. For a curve, the magnification is not constant -- it depends on where
you are -- and that forces one new idea: \emph{control the wild factor first}.

\begin{example}[A quadratic limit]
\label{ch06:ex-quad}
Prove that $\displaystyle \lim_{x\to 3} x^2 = 9$.
\end{example}

\textbf{Scratch work.} The output gap is
\[
  \abs{x^2 - 9} = \abs{(x-3)(x+3)} = \abs{x-3}\,\abs{x+3}.
\]
One factor, $\abs{x-3}$, is the input gap---controlled directly via
$\delta$. The other factor, $\abs{x+3}$, varies: near $x=3$ it is
about $6$, but it is not constant, so we must establish a bound on it.
The solution is a standard two-step technique.

\emph{Step one: bound the variable factor with a preliminary clamp.} Decide in advance
that $\delta$ will be \emph{at most} $1$. If $\abs{x-3}<1$, then $2<x<4$, so
\[
  5 < x+3 < 7, \quad\text{hence}\quad \abs{x+3} < 7.
\]
As long as we keep $x$ within $1$ of $3$, the factor $\abs{x+3}$ cannot exceed $7$.
(Any valid upper bound works; $7$ is convenient.)

\emph{Step two: use the bound to size $\delta$.} Now, whenever $\abs{x-3}<1$,
\[
  \abs{x^2-9} = \abs{x-3}\,\abs{x+3} < 7\,\abs{x-3}.
\]
To force this under $\eps$ it suffices to have $\abs{x-3} < \eps/7$. We need
\emph{both} clamps at once -- $\abs{x-3}<1$ to validate the bound, and
$\abs{x-3}<\eps/7$ to land in the band -- so we take the smaller of the two:
\[
  \delta = \min\!\left(1,\ \frac{\eps}{7}\right).
\]

\textbf{Proof.} Let $\eps>0$ be given. Choose $\delta = \min(1,\eps/7)$. Suppose
$0<\abs{x-3}<\delta$. Since $\delta\le 1$, we have $\abs{x-3}<1$, so $2<x<4$ and
therefore $\abs{x+3}<7$. Since also $\delta\le\eps/7$, we have $\abs{x-3}<\eps/7$.
Combining,
\[
  \abs{x^2 - 9} = \abs{x-3}\,\abs{x+3} < \frac{\eps}{7}\cdot 7 = \eps.
\]
Hence $\lim_{x\to 3} x^2 = 9$. $\qquad\blacksquare$

\begin{intuitionbox}
The $\min$ is not a trick; it is the natural answer to a genuine two-part demand.
One part of the requirement (``keep the slope factor bounded'') sets a ceiling on
$\delta$; the other part (``land inside the $\eps$ band'') sets another. To
satisfy both promises simultaneously you take the tighter one. Whenever you see
$\delta = \min(\dots,\dots)$ in an epsilon-delta proof, read it as ``I made two
promises and I am keeping both.'' Figure~\ref{ch06:fig-quad-clamp} illustrates
how the two promises divide the work.
\end{intuitionbox}

\begin{figure}[htbp]
\centering
\begin{tikzpicture}
\begin{axis}[
  width=10.5cm, height=6.8cm,
  title style={font=\small},
  title={Taming the wild factor: the two-step clamp for $f(x)=x^2$ near $a=3$},
  axis lines=middle,
  xlabel={$x$}, ylabel={$y$},
  xmin=1.4, xmax=4.6, ymin=2.0, ymax=18.0,
  xtick={2, 3, 4},
  xticklabels={$2$, $a=3$, $4$},
  ytick={9},
  yticklabels={$L=9$},
  clip=false,
]
% Shaded regions first
\addplot[draw=none, fill=gray!12] coordinates {(2,2)(4,2)(4,18)(2,18)};
\addplot[draw=none, fill=green!16] coordinates {(1.4,10.4)(4.6,10.4)(4.6,7.6)(1.4,7.6)};
\addplot[draw=none, fill=blue!14] coordinates {(2.8,2)(3.2,2)(3.2,18)(2.8,18)};

% Boundary lines
\draw[green!50!black, dashed] (axis cs:1.4,10.4) -- (axis cs:4.6,10.4);
\draw[green!50!black, dashed] (axis cs:1.4,7.6) -- (axis cs:4.6,7.6);
\draw[blue!50!black, dashed] (axis cs:2.8,2) -- (axis cs:2.8,18);
\draw[blue!50!black, dashed] (axis cs:3.2,2) -- (axis cs:3.2,18);

% Function curve and target
\addplot[curveblue, very thick, domain=1.5:4.2, samples=60] {x^2};
\draw[gray, dashed] (axis cs:3,2) -- (axis cs:3,9);
\draw[gray, dashed] (axis cs:1.4,9) -- (axis cs:3,9);
\fill[white] (axis cs:3,9) circle (2.4pt);
\draw[curveblue, thick] (axis cs:3,9) circle (2.4pt);

% Annotations in clear space
\node[gray!60!black, font=\scriptsize, anchor=west] at (axis cs:2.05,14.5) {clamp: $\abs{x-3}<1$};
\node[green!50!black, font=\scriptsize, anchor=east] at (axis cs:1.35,10.4) {$L+\eps$};
\node[green!50!black, font=\scriptsize, anchor=east] at (axis cs:1.35,7.6) {$L-\eps$};
\node[curveblue, anchor=south east] at (axis cs:4.1,16.8) {$y=x^2$};
\node[blue!50!black, font=\scriptsize, anchor=south west] at (axis cs:3.05,2.4) {$\delta = \eps/7$};
\end{axis}
\end{tikzpicture}
\caption{The two-step clamp for $f(x)=x^2$ near $a=3$. Step~1 sets a preliminary
ceiling $\delta\le 1$, confining $x$ to the wide gray region $(2,4)$ where the
slope factor $\abs{x+3}$ can never exceed $7$. Step~2 uses this ceiling to size
the input strip $\delta = \eps/7$ (the narrow blue strip). Inside this strip,
the curve $y=x^2$ stays trapped between $L-\eps$ and $L+\eps$ (the green band).}
\label{ch06:fig-quad-clamp}
\end{figure}

\begin{whybox}
Why does the curve need a $\min$ whereas the linear function did not? For a linear
function, the magnification factor is constant, so a single choice of $\delta$ suffices.
For $x^2$, the factor varies with $x$. We therefore restrict $x$ to a neighborhood
where the factor is bounded by a constant, and then choose $\delta$ small enough to
satisfy the target tolerance. The underlying principle is:
\emph{a controllable input gap multiplied by a bounded factor yields a controllable output gap}.
Every product limit argument, including the product limit law below, relies on this
principle that a vanishing factor multiplied by a bounded factor vanishes.
\end{whybox}

\begin{trapbox}
\textbf{Do not forget to bound the variable factor.} A common wrong ``proof'' of
$\lim_{x\to3}x^2=9$ writes $\abs{x^2-9}=\abs{x-3}\abs{x+3}$ and then sets
$\delta=\eps/\abs{x+3}$. But $\delta$ is not allowed to depend on $x$ -- $x$ is
the variable we are quantifying over, not a fixed number you may plug into your
choice of $\delta$. The Prover must commit to a single $\delta$ depending only on
$\eps$ (and on the fixed data $a$), \emph{before} any particular $x$ is examined.
The preliminary clamp $\abs{x-3}<1$ exists precisely so that we can replace the
forbidden, $x$-dependent $\abs{x+3}$ by the fixed constant $7$.
\end{trapbox}

\begin{labbox}
Prove $\displaystyle\lim_{x\to 2}x^2 = 4$ on your own. Clamp $\abs{x-2}<1$ so
that $1<x<3$ and $\abs{x+2}<5$, then take $\delta=\min(1,\eps/5)$. Compare with
the $x=3$ case: the amplifier ceiling changed from $7$ to $5$ because you are
sitting at a less steep part of the parabola. The $\delta$ you are allowed to use
literally tracks the steepness of the curve at the point -- another preview of
the derivative.
\end{labbox}

\begin{recoverybox}
\textbf{Pause and think.} What would happen if we used a preliminary clamp of
$\delta \le 10$ instead of $\delta \le 1$ when proving $\lim_{x\to 3} x^2 = 9$?
If $\abs{x-3} < 10$, then $-7 < x < 13$, so $\abs{x+3} < 16$. The proof
still goes through: we simply choose $\delta = \min(10,\, \eps/16)$.
The initial clamp of $1$ is an arbitrary convenient choice---any positive radius that
bounds the variable factor away from infinity serves the same purpose.
\end{recoverybox}

Let us record the general lesson as a small lemma, because it is the engine
of all the later laws.

\begin{lemma}[Small times bounded is small]
\label{ch06:lem-stb}
Suppose that on some punctured interval around $a$ we can write $f(x)-L =
g(x)\,h(x)$, where $\abs{g(x)} \le M$ for some constant $M>0$ on that interval,
and $\lim_{x\to a} h(x) = 0$. Then $\lim_{x\to a} f(x) = L$.
\end{lemma}

\begin{proof}
Let $\eps>0$. Since $\lim_{x\to a}h(x)=0$, there is $\delta_1>0$ with
$0<\abs{x-a}<\delta_1 \Rightarrow \abs{h(x)} < \eps/M$. Let $\delta_2>0$ be a
radius small enough that the bound $\abs{g(x)}\le M$ holds for
$0<\abs{x-a}<\delta_2$. Put $\delta=\min(\delta_1,\delta_2)$. Then for
$0<\abs{x-a}<\delta$,
\[
  \abs{f(x)-L} = \abs{g(x)}\,\abs{h(x)} < M\cdot\frac{\eps}{M} = \eps. \qedhere
\]
\end{proof}

The quadratic proof was exactly this lemma with $g(x)=x+3$ (bounded by $M=7$ near
$3$) and $h(x)=x-3$ (heading to $0$). Hold this pattern; it returns.

\section{The limit laws, and why they are theorems}
\label{ch06:laws}

In practice nobody computes $\lim_{x\to 3}(x^2+5x)$ from scratch with epsilon and
delta. Instead they say: the limit of a sum is the sum of the limits, the limit
of a product is the product of the limits, and so on -- so the answer is $9+15 =
24$. These \emph{limit laws} are what make limits usable. But they are not
axioms. Each one is a \emph{theorem} that must be \emph{earned} from
Definition~\ref{ch06:def-limit}. The intuition (errors add, small times bounded
is small) tells you the laws should be true; the epsilon-delta machine is what
\emph{makes} them true. Let us prove the two central ones and see the intuition
become rigor.

Assume throughout that $\lim_{x\to a} f(x) = L$ and $\lim_{x\to a} g(x) = M$.

\begin{theorem}[Sum law]
\label{ch06:thm-sum}
$\displaystyle \lim_{x\to a}\big(f(x)+g(x)\big) = L + M.$
\end{theorem}

\begin{intuitionbox}
The idea before the proof. The total error of the sum, $\abs{(f+g)-(L+M)}$, is at
most the error of $f$ plus the error of $g$ (by the triangle inequality, which
states that the direct distance $\abs{p+q}$ is never greater than the detour
$\abs{p}+\abs{q}$). If you want the \emph{combined} error under $\eps$, just
demand each piece be under $\eps/2$. Halve the tolerance, hand half to each
function, collect their separate $\delta$'s, and use the tighter one. ``Split
the budget'' is the entire strategy.
\end{intuitionbox}

\begin{proof}
Let $\eps>0$. Apply the hypotheses with the tolerance $\eps/2$:
\begin{itemize}
\item Since $\lim_{x\to a}f=L$, there is $\delta_1>0$ with
  $0<\abs{x-a}<\delta_1 \Rightarrow \abs{f(x)-L}<\eps/2$.
\item Since $\lim_{x\to a}g=M$, there is $\delta_2>0$ with
  $0<\abs{x-a}<\delta_2 \Rightarrow \abs{g(x)-M}<\eps/2$.
\end{itemize}
Let $\delta=\min(\delta_1,\delta_2)$. For $0<\abs{x-a}<\delta$ \emph{both} bounds
hold at once, so by the triangle inequality
\[
  \abs{\big(f(x)+g(x)\big)-(L+M)}
  = \abs{\big(f(x)-L\big)+\big(g(x)-M\big)}
  \le \abs{f(x)-L}+\abs{g(x)-M}
  < \frac{\eps}{2}+\frac{\eps}{2} = \eps.
\]
Therefore $\lim_{x\to a}(f+g)=L+M$.
\end{proof}

\begin{trapbox}
\textbf{Why $\eps/2$ and not $\eps$?} If you only forced each piece under $\eps$,
the triangle inequality would give you $\abs{(f+g)-(L+M)} < \eps+\eps = 2\eps$ --
twice the target, a miss. Splitting the budget in half is what makes the two
halves \emph{add up} to exactly $\eps$. The choice $\eps/2$ is reverse-engineered
from the triangle inequality's ``$+$''. This budget-splitting is one of the most
reused tactics in all of analysis; it will reappear as ``$\eps/3$ tricks'' and
``$\eps/n$ tricks'' whenever finitely many errors must be controlled together.
\end{trapbox}

\begin{theorem}[Product law]
\label{ch06:thm-prod}
$\displaystyle \lim_{x\to a}\big(f(x)\,g(x)\big) = L\,M.$
\end{theorem}

\begin{intuitionbox}
Products are subtler than sums because the errors interact. Write $f = L + u$ and
$g = M + v$, where $u=f-L$ and $v=g-M$ are the (small) errors. Then
\[
  fg = (L+u)(M+v) = LM + Lv + Mu + uv,
\]
so the product's error is
\[
  fg - LM = \underbrace{Lv}_{\text{const}\times\text{small}}
  + \underbrace{Mu}_{\text{const}\times\text{small}}
  + \underbrace{uv}_{\text{small}\times\text{small}}.
\]
Every one of the three pieces is small: the first two are a fixed constant times
something heading to $0$, and the third is even smaller, a product of two
vanishing errors. Make each of the three under $\eps/3$ and add. That is the
proof; the rest is bookkeeping to turn ``small'' into actual inequalities.
\end{intuitionbox}

\begin{proof}
Let $\eps>0$. We control the three pieces $\abs{Lv},\ \abs{Mu},\ \abs{uv}$, where
$u=f(x)-L$ and $v=g(x)-M$, aiming each below $\eps/3$.

\emph{The cross term $\abs{uv}$.} Pick $\delta_1,\delta_2>0$ so that
$\abs{u}=\abs{f(x)-L}<\sqrt{\eps/3}$ for $0<\abs{x-a}<\delta_1$ and
$\abs{v}=\abs{g(x)-M}<\sqrt{\eps/3}$ for $0<\abs{x-a}<\delta_2$. Then
$\abs{uv}<\eps/3$.

\emph{The term $\abs{Mu}$.} Pick $\delta_3>0$ so that $\abs{u}<\dfrac{\eps}{3(\abs{M}+1)}$
for $0<\abs{x-a}<\delta_3$. Then $\abs{Mu}\le\abs{M}\,\abs{u} \le
\dfrac{\abs{M}}{\abs{M}+1}\cdot\dfrac{\eps}{3} < \dfrac{\eps}{3}$. (We divide by
$\abs{M}+1$ rather than $\abs{M}$ so the bound is safe even when $M=0$.)

\emph{The term $\abs{Lv}$.} Symmetrically, pick $\delta_4>0$ so that
$\abs{v}<\dfrac{\eps}{3(\abs{L}+1)}$ for $0<\abs{x-a}<\delta_4$, giving
$\abs{Lv}<\eps/3$.

Now set $\delta=\min(\delta_1,\delta_2,\delta_3,\delta_4)$. For
$0<\abs{x-a}<\delta$ all four conditions hold simultaneously, so by the triangle
inequality
\[
  \abs{f(x)g(x)-LM} = \abs{Lv + Mu + uv}
  \le \abs{Lv}+\abs{Mu}+\abs{uv}
  < \frac{\eps}{3}+\frac{\eps}{3}+\frac{\eps}{3} = \eps.
\]
Hence $\lim_{x\to a}(fg)=LM$.
\end{proof}

\begin{remark}
\label{ch06:rem-prod}
The proof above is the direct, self-contained one, and it shows clearly where each
term goes. A slicker route uses Lemma~\ref{ch06:lem-stb}: since $f$ has a limit it
is bounded near $a$, and $fg-LM = f\,(g-M) + M\,(f-L)$ is then ``bounded times
small plus constant times small,'' each piece $\to 0$. Both proofs are really the
same observation -- \emph{small errors, multiplied by quantities that stay
bounded, stay small} -- which is why we isolated that lemma earlier.
\end{remark}

\begin{aestheticbox}
The sum law and product law are both consequences of two basic properties of distance:
the triangle inequality ($\abs{p+q}\le\abs{p}+\abs{q}$), and the principle that
a vanishing quantity multiplied by a bounded quantity vanishes.
The entire arithmetic of limits rests on these foundations:
simple geometric properties of distance and bounds rather than arbitrary algebraic rules.
\end{aestheticbox}

For completeness, here are the remaining standard laws. The constant-multiple and
difference laws are immediate from the two above (take $g$ constant, or replace
$g$ by $-g$). The quotient law needs one extra idea -- keeping the denominator
away from zero -- which we state and sketch.

\begin{theorem}[The limit laws]
\label{ch06:thm-laws}
If $\lim_{x\to a}f(x)=L$ and $\lim_{x\to a}g(x)=M$, then:
\begin{enumerate}
\item (Constant multiple) $\displaystyle\lim_{x\to a} c\,f(x) = cL$ for any
  constant $c$.
\item (Sum/difference) $\displaystyle\lim_{x\to a}\big(f(x)\pm g(x)\big) = L\pm M$.
\item (Product) $\displaystyle\lim_{x\to a} f(x)g(x) = LM$.
\item (Quotient) $\displaystyle\lim_{x\to a}\frac{f(x)}{g(x)} = \frac{L}{M}$,
  provided $M\ne 0$.
\end{enumerate}
\end{theorem}

\begin{proof}[Sketch of the quotient law]
It is enough, given the product law, to handle $\lim_{x\to a}\frac{1}{g(x)} =
\frac{1}{M}$ when $M\ne0$; then $\frac{f}{g}=f\cdot\frac1g$ follows. Compute the
error:
\[
  \abs{\frac{1}{g(x)} - \frac{1}{M}}
  = \frac{\abs{M - g(x)}}{\abs{g(x)}\,\abs{M}}.
\]
The numerator $\abs{M-g(x)}\to0$ -- that is small. The danger is the
$\abs{g(x)}$ in the denominator: if it could be near $0$, the fraction could blow
up. So we first secure the denominator. Since $M\ne0$, apply the limit of $g$
with tolerance $\abs{M}/2$: there is a radius within which $\abs{g(x)-M} <
\abs{M}/2$, which implies $\abs{g(x)} \ge \abs{M} - \abs{g(x)-M} > \abs{M}/2 > 0$
by the reverse triangle inequality. On that region the fraction is at most
$\frac{2}{\abs{M}^2}\,\abs{M-g(x)}$ -- once again ``small times bounded'' -- and
shrinking the radius further forces it under any $\eps$. The structure is
identical to the quadratic proof: clamp first to control the dangerous factor,
then shrink to hit the band.
\end{proof}

\begin{connectionbox}
The limit laws say the map ``$f \mapsto \lim_{x\to a}f$'' respects addition and
multiplication: it sends sums to sums and products to products. In the language
of later mathematics, taking the limit at $a$ is a \emph{homomorphism} -- a
structure-preserving dictionary that carries operations in one algebraic system
into matching operations in another -- from the algebra of functions having a
limit at $a$ into $\R$. The same structural fact --
``the operation is linear and multiplicative'' -- recurs for the derivative
(linear, with a product rule) and the integral (linear). Spotting which
operations preserve which structures is one of the organizing instincts of all
higher mathematics, and limits are where it starts.
\end{connectionbox}

\section{When the game cannot be won: limits that do not exist}
\label{ch06:no-limit}

The definition's real power is that it can also certify \emph{failure}. To show a
limit does \emph{not} exist, we play the Skeptic: we exhibit one tolerance $\eps$
that defeats every possible $\delta$.

\begin{example}[A jump]
\label{ch06:ex-jump}
The Heaviside step $H(x)$ -- a simple switch function defined by $H(x)=0$ for
$x<0$ and $H(x)=1$ for $x\ge 0$ -- has no limit at $0$.
\end{example}

\textbf{Argument.} Suppose, for contradiction, some $L$ were the limit. Take the
Skeptic's tolerance $\eps=\tfrac12$. For \emph{any} $\delta>0$, the punctured
interval $(-\delta,\delta)\setminus\{0\}$ contains both negative points (where
$H=0$) and positive points (where $H=1$). For the promise to hold we would need
both $\abs{0-L}<\tfrac12$ and $\abs{1-L}<\tfrac12$, i.e. $L$ within $\tfrac12$ of
\emph{both} $0$ and $1$. But $0$ and $1$ are a distance $1$ apart, and no single
point is within $\tfrac12$ of both. So $\eps=\tfrac12$ defeats every $\delta$;
the limit cannot exist. $\qquad\blacksquare$

\begin{figure}[htbp]
\centering
\begin{tikzpicture}
\begin{axis}[
  width=10.5cm, height=5.5cm,
  title style={font=\small},
  title={Why a jump fails the $\eps$--$\delta$ game (tested with candidate $L=1$ and $\eps=0.4$)},
  axis lines=middle, xlabel={$x$}, ylabel={$y$},
  xlabel style={anchor=south west},
  ylabel style={anchor=south east},
  xmin=-1.6, xmax=1.6, ymin=-0.4, ymax=1.6,
  xtick={-0.6,0.6}, xticklabels={$-\delta$,$+\delta$},
  xticklabel style={fill=white, fill opacity=0.9, text opacity=1, inner sep=1.5pt},
  ytick={0.6,1,1.4}, yticklabels={$L-\eps$,$L$,$L+\eps$},
  yticklabel style={anchor=east, fill=white, fill opacity=0.9, text opacity=1, inner sep=1.5pt},
  axis on top=true,
  clip=false,
]
\addplot[draw=none, fill=green!15] coordinates {(-1.6,1.4)(1.6,1.4)(1.6,0.6)(-1.6,0.6)};
\addplot[draw=none, fill=blue!10] coordinates {(-0.6,-0.4)(0.6,-0.4)(0.6,1.6)(-0.6,1.6)};
\addplot[tangentred, line width=1.5pt, domain=-0.6:-0.02, samples=2] {0};
\addplot[curveblue, very thick, domain=-1.6:-0.6, samples=2] {0};
\addplot[curveblue, very thick, domain=0.02:1.6, samples=2] {1};
\draw[gray, dashed] (axis cs:-0.6,-0.4) -- (axis cs:-0.6,1.6);
\draw[gray, dashed] (axis cs:0.6,-0.4) -- (axis cs:0.6,1.6);
\fill[white] (axis cs:0,0) circle (2.2pt);
\draw[curveblue] (axis cs:0,0) circle (2.2pt);
\fill[curveblue] (axis cs:0,1) circle (2.2pt);
\node[tangentred, font=\scriptsize, anchor=south] at (axis cs:-0.8, 0.08) {escapes band: $H(x)=0$};
\end{axis}
\end{tikzpicture}
\caption{Negation of the limit at a jump. Testing candidate $L=1$ with $\eps=0.4$ creates a horizontal band $[0.6,1.4]$. Any proposed vertical strip $(-\delta,\delta)$ inevitably contains negative inputs (red segment) where $H(x)=0$, lying entirely outside the band. No $\delta$ can ever succeed.}
\label{ch06:fig-jump-eps}
\end{figure}

\begin{example}[Infinite oscillation]
\label{ch06:ex-sin}
$\displaystyle\lim_{x\to 0}\sin\!\frac{1}{x}$ does not exist.
\end{example}

As $x\to0$, the input $1/x$ races off to infinity, so $\sin(1/x)$ runs through
its full range $[-1,1]$ infinitely often -- it never settles. No matter how thin
a strip $(-\delta,\delta)$ the Prover proposes, that strip still contains inputs
where $\sin(1/x)=1$ and inputs where $\sin(1/x)=-1$ (because $1/x$ sweeps through
infinitely many periods). With the Skeptic's tolerance $\eps=1$, no candidate $L$
can be within $1$ of both $+1$ and $-1$ at once. Every $\delta$ leaks; the limit
fails to exist. Figure~\ref{ch06:fig-sin} shows the graph refusing to settle.

\begin{figure}[htbp]
\centering
\begin{tikzpicture}
\begin{axis}[
  width=12cm, height=5.4cm,
  title style={font=\small},
  title={Failure to settle: infinite oscillation on $[-1, 1]$},
  axis lines=middle, xlabel={$x$}, ylabel={$y$},
  xlabel style={anchor=south west},
  ylabel style={anchor=south east},
  xmin=-0.38, xmax=0.38, ymin=-1.4, ymax=1.4,
  ytick={-1,1}, xtick={0},
  yticklabel style={fill=white, fill opacity=0.9, text opacity=1, inner sep=1.5pt},
]
\addplot[curveblue, thick, samples=200, domain=0.004:0.35] {sin(deg(1/x))};
\addplot[curveblue, thick, samples=200, domain=-0.35:-0.004] {sin(deg(1/x))};
\addplot[green!50!black, dashed, domain=-0.38:0.38, samples=2] {1};
\addplot[green!50!black, dashed, domain=-0.38:0.38, samples=2] {-1};
\end{axis}
\end{tikzpicture}
\caption{$y=\sin(1/x)$ near $0$. The oscillations bunch up infinitely as $x\to0$,
forever sweeping the whole band from $-1$ to $1$. No vertical strip around $0$,
however thin, can trap the graph inside any band thinner than the full height
$2$, so no limit exists. (Pixels cannot draw the infinitely many wiggles closest
to $0$; imagine them packed ever tighter.)}
\label{ch06:fig-sin}
\end{figure}

\begin{connectionbox}
The jump in $H$ fails for a ``two-sided disagreement'' reason: the left approach
wants $L=0$, the right wants $L=1$. This motivates \emph{one-sided limits} --
limits that examine function values when approaching strictly from the left
($x\to a^-$) or strictly from the right ($x\to a^+$) -- where the input strip is
simply replaced by a half-interval. A two-sided limit exists exactly when both
one-sided limits exist and agree -- which they do not for $H$. We develop
one-sided limits and the closely related idea of continuity in the chapters just
ahead; here, note only that the epsilon-delta frame absorbs them with a single
change (replace the symmetric strip by a half-strip).
\end{connectionbox}

\section{Making the intuition and the formalism agree}
\label{ch06:agree}

We end where the chapter began: with the relationship between the picture and the
sentence. They had better agree -- a definition that contradicted the clear cases
would be a bad definition, not a brave one. Check the handshake from both
sides.

\textbf{Formalism honors intuition.} Every limit you would compute ``by
inspection'' -- $\lim_{x\to2}(3x+1)=7$, $\lim_{x\to3}x^2=9$, $\lim_{x\to a}c=c$ --
comes out the expected way under Definition~\ref{ch06:def-limit}, as we proved or
sketched. The definition does not overturn your instincts for the easy cases; it
ratifies them. What it adds is a referee for the \emph{hard} cases, where
instinct stalls: removable holes like $\frac{x^2-1}{x-1}$, jumps like $H$,
oscillations like $\sin(1/x)$. There the picture alone gives no confident answer,
and the definition decides.

\textbf{Intuition tames the formalism.} Conversely, you should never run the
epsilon-delta machine blind. Before writing a single $\delta$, look at the
picture and \emph{predict} the answer; the formal proof is then a verification,
not a search in the dark. And the scratch-work phase -- working backward from
$\abs{f(x)-L}<\eps$ to find $\delta$ -- is just the picture turned into algebra:
``how tight a strip keeps the graph in this band?'' The two are one activity in
two notations.

\begin{happeningbox}
Here is the cleanest way to hold the whole chapter in one thought. The
intuitive limit and the epsilon-delta limit are not two different limits that
happen to agree; they are the \emph{same} statement told at two resolutions. ``$f(x)$
approaches $L$'' is the statement at human resolution -- fast, visual,
occasionally wrong at the edges. ``$\forall\eps\,\exists\delta\dots$'' is the
identical statement at full resolution -- slow, symbolic, never wrong. You learn
both because you think in the first and prove in the second, and a working
mathematician fluidly translates between them mid-sentence. When that
translation becomes automatic, the epsilon-delta definition stops being an
obstacle and becomes what it truly is: the exact written form of something you
already see.
\end{happeningbox}

\begin{walkbox}
A question to turn over on a walk. We insisted the order of quantifiers is
``$\forall\eps\,\exists\delta$.'' Imagine a stricter notion in which a single
$\delta$ must work for \emph{all} $\eps$ simultaneously: $\exists\delta\,
\forall\eps$. Which functions, if any, satisfy \emph{that}? (Hint: it would force
$\abs{f(x)-L}<\eps$ for every $\eps>0$ on a fixed punctured interval, hence
$f(x)=L$ there.) You will discover that the strict version describes only
locally-constant functions -- which is exactly why the ordinary definition uses
the looser, correct order. The grammar is doing real work.
\end{walkbox}

\subsection{From verification to calculation: looking ahead}

The epsilon--delta definition is an instrument of verification, not calculation: it checks whether
a suspected limit $L$ is correct, but it does not tell you how to find $L$. Having proved the
limit laws, we can now compute limits without building custom $\eps$--$\delta$ proofs for each one.
In Chapter~7, we turn to algebraic techniques---factoring, conjugates, and the squeeze theorem---for
evaluating limits directly.

\section{Exercises}
\label{ch06:exercises}

\begin{exercise}[Mental checks: reading the quantifiers]
Answer these three checks without writing an algebraic proof:
\begin{enumerate}
\item If a function satisfies $\abs{f(x)-3} < 0.05$ whenever $0 < \abs{x-1} < 0.01$,
  did the Prover just successfully answer a challenge for $\eps=0.05$ or
  $\eps=0.01$?
\item In an epsilon--delta argument, why is the choice $\delta = \eps / \abs{x-a}$
  strictly illegal?
\item If $f(2) = 500$ but $\lim_{x\to 2} f(x) = 7$, does the value $500$ prevent
  the Prover from finding a winning $\delta$ for $\eps = 0.1$? Explain in one
  sentence.
\end{enumerate}
\end{exercise}

\begin{exercise}[Linear, by the template]
Using the two-step template (scratch work, then proof), prove from
Definition~\ref{ch06:def-limit} that $\displaystyle\lim_{x\to 4}(2x-3)=5$.
Identify the rule $\eps\mapsto\delta$ explicitly and relate the constant in it to
the slope of the line.
\end{exercise}

\begin{exercise}[Negative slope]
Prove that $\displaystyle\lim_{x\to 1}(8-5x)=3$. Be careful that a negative slope
does not change the size of the magnification: explain why $\delta=\eps/5$ (not
$\eps/(-5)$) is the right response, and where the absolute value enters.
\end{exercise}

\begin{exercise}[Quadratic with a clamp]
Prove that $\displaystyle\lim_{x\to 4}x^2=16$. State your preliminary clamp,
derive the resulting ceiling on $\abs{x+4}$, and write
$\delta=\min(1,\eps/9)$ (or justify a different constant if you clamp
differently).
\end{exercise}

\begin{exercise}[A reciprocal]
Prove that $\displaystyle\lim_{x\to 2}\frac{1}{x}=\frac12$ directly from the
definition. (Compute $\abs{\frac1x-\frac12}=\frac{\abs{x-2}}{2\abs{x}}$, clamp
$\abs{x-2}<1$ so that $1<x<3$ and hence $\abs{x}>1$, then choose
$\delta=\min(1,2\eps)$.) Point out exactly where you used the clamp to keep the
denominator away from $0$.
\end{exercise}

\begin{exercise}[Negating the definition]
Write out, with quantifiers in the correct order, the precise statement
``$\lim_{x\to a}f(x)\ne L$.'' Then use it to prove that
$\displaystyle\lim_{x\to 0}\frac{\abs{x}}{x}$ does not exist. (Take $\eps=1$ and
show every $\delta$ contains points where the function is $+1$ and points where
it is $-1$.)
\end{exercise}

\begin{exercise}[The $\eps/2$ budget, examined]
In the proof of the sum law we chose tolerance $\eps/2$ for each function.
\textbf{(a)} Redo the final inequality using $\eps/3$ for each piece and confirm
it still works (you land at $2\eps/3<\eps$). \textbf{(b)} Show that using $\eps$
for each piece \emph{fails} to prove the claim, and pinpoint the exact step where
the argument breaks. \textbf{(c)} Conclude that any split summing to at most
$\eps$ works, and that $\eps/2$ is simply the tidiest.
\end{exercise}

\begin{exercise}[Constant multiple law]
Prove from the definition that $\displaystyle\lim_{x\to a}c\,f(x)=cL$ whenever
$\lim_{x\to a}f(x)=L$. Treat $c=0$ separately and explain why that case needs no
hypothesis on $f$ at all.
\end{exercise}

\begin{exercise}[Small times bounded, applied]
Use Lemma~\ref{ch06:lem-stb} to prove $\displaystyle\lim_{x\to0}x\sin\frac1x=0$
even though $\sin\frac1x$ has no limit at $0$. Identify $g$ and $h$ in the lemma,
and the bound $M$. (This is a striking case: a wildly oscillating factor,
multiplied by something heading to $0$, is tamed completely.)
\end{exercise}

\begin{exercise}[A removable hole]
Let $\displaystyle f(x)=\frac{x^2-4}{x-2}$ for $x\ne2$ and $f(2)=100$. Prove that
$\lim_{x\to2}f(x)=4$ directly from the definition, and explain in one sentence
why the absurd value $f(2)=100$ is invisible to the limit. Which inequality in
Definition~\ref{ch06:def-limit} is responsible for that invisibility?
\end{exercise}

\begin{exercise}[Conceptual: order of quantifiers, proof]
Prove that if a single $\delta>0$ satisfied the condition
``$0<\abs{x-a}<\delta\Rightarrow\abs{f(x)-L}<\eps$'' for \emph{every} $\eps>0$
at once, then $f(x)=L$ for all $x$ in the punctured interval $0<\abs{x-a}<\delta$.
Conclude that the order $\forall\eps\,\exists\delta$ in the definition is
essential and cannot be reversed without trivializing it.
\end{exercise}

\begin{exercise}[Proof: uniqueness of the limit]
Prove that a limit, if it exists, is unique: if $\lim_{x\to a}f(x)=L$ and
$\lim_{x\to a}f(x)=L'$, then $L=L'$. (Suppose $L\ne L'$ and let
$\eps=\abs{L-L'}/2>0$. Use the definition for each limit to trap $f(x)$ within
$\eps$ of both $L$ and $L'$ simultaneously on a small punctured interval, and
derive a contradiction from the triangle inequality.) This is the theorem that
licenses the phrase ``\emph{the} limit.''
\end{exercise}
